% ============================================================================
%  SISEPUEDE Energy Calibration — Country Report (template)
% ----------------------------------------------------------------------------
%  Location: sisepuede/calibration/docs/energy_calibration_report_template.tex
%
%  Fill in the five \newcommand{...} values below (country ISO, country
%  name, target year, and the two output directories produced by
%  `energy_calibration()`), then compile with pdflatex. The directory
%  paths are passed straight to \graphicspath / \input, so absolute paths
%  always work; relative paths are resolved against the directory where
%  pdflatex is invoked (not where this template lives), so prefer
%  absolute paths unless you compile from a known location.
%
%  Plots and tables generated by the pipeline are picked up automatically
%  as long as the file names follow the pipeline's convention:
%
%      <plotsDir>/<ISO>_baseline_discrepancy_ratio.png
%      <plotsDir>/<ISO>_baseline_discrepancy_rel_error.png
%      <plotsDir>/<ISO>_before_after_discrepancy_ratio.png
%      <plotsDir>/<ISO>_before_after_discrepancy_rel_error.png
%      <plotsDir>/<ISO>_primary_timeseries.png
%      <plotsDir>/<ISO>_primary_bar.png
%      <plotsDir>/<ISO>_fuelmix_<SECTOR>_timeseries.png
%      <plotsDir>/<ISO>_fuelmix_<SECTOR>_bar.png
%      <tablesDir>/<ISO>_values_table.tex
%      <tablesDir>/<ISO>_improvement_table.tex
%      <tablesDir>/knobs/knobs_<BALANCE>_<PRODUCT>.tex
%
%  For the per-sector fuel-mix and per-pair knob blocks, duplicate the
%  marked example block once per sector / (balance, product) pair.
% ============================================================================

\documentclass[10pt]{article}

\usepackage[margin=1in]{geometry}
\usepackage{graphicx}
\usepackage{booktabs}
\usepackage[table]{xcolor}
\usepackage{float}
\usepackage{caption}
\usepackage{subcaption}
\usepackage{hyperref}

% ── User-supplied values ────────────────────────────────────────────────────
\newcommand{\countryISO}{PER}                                                  % e.g. PER, MEX, COL
\newcommand{\countryName}{Peru}                                                % human-readable name
\newcommand{\targetYear}{2018}                                                 % calibration target year
\newcommand{\plotsDir}{/Users/dianamendez/sisepuede-data/plots}                % path to plots/
\newcommand{\tablesDir}{/Users/dianamendez/sisepuede-data/tables}              % path to tables/
\newcommand{\runTag}{}                                                         % optional run tag (e.g. gamma100); leave empty for none

% ── Path helpers (do not edit) ──────────────────────────────────────────────
%   When \runTag is non-empty, filenames get an extra "_<runTag>" suffix
%   and the per-pair knob tables live in "knobs_<runTag>/", matching what
%   energy_calibration() writes when called with a tag.
\graphicspath{{\plotsDir/}}

% Internal: tag suffix (expands to "" if \runTag is empty, "_<tag>" otherwise)
\makeatletter
\newcommand{\tagSuffix}{\ifx\runTag\@empty\else_\runTag\fi}
\makeatother

\newcommand{\plotFile}[1]{\plotsDir/\countryISO_#1\tagSuffix.png}
\newcommand{\tableFile}[1]{\tablesDir/\countryISO_#1\tagSuffix.tex}
\newcommand{\knobFile}[2]{\tablesDir/knobs\tagSuffix/knobs_#1_#2}

% ── Title ───────────────────────────────────────────────────────────────────
\title{SISEPUEDE Energy Calibration Report \\ \large \countryName\ (\countryISO) --- target year \targetYear}
\author{}
\date{}

\begin{document}
\maketitle
\tableofcontents
\bigskip


% ============================================================================
\section{Overview}
% ============================================================================

This report summarises the SISEPUEDE energy calibration run for
\countryName\ (\countryISO) with calibration target year \targetYear.
The calibration adjusts SISEPUEDE energy inputs so that the model's
outputs, aggregated through the IEA crosswalk, match the IEA World
Energy Balances at the target year. The two-phase procedure scales
sector totals first (Phase~1) and then re-balances the fuel mix within
each sector (Phase~2, simplex-constrained).


% ============================================================================
\section{Baseline discrepancy (before calibration)}
% ============================================================================

The figures below show the SISEPUEDE/IEA discrepancy at the target
year prior to calibration, as a ratio (SSP/IEA, reference~$=1$) and as
a relative error ((SSP$-$IEA)/IEA, reference~$=0$).

\begin{figure}[H]
    \centering
    \begin{subfigure}{0.49\textwidth}
        \includegraphics[width=\linewidth]{\plotFile{baseline_discrepancy_ratio}}
        \caption{Ratio SSP/IEA.}
    \end{subfigure}\hfill
    \begin{subfigure}{0.49\textwidth}
        \includegraphics[width=\linewidth]{\plotFile{baseline_discrepancy_rel_error}}
        \caption{Relative error.}
    \end{subfigure}
    \caption{Baseline discrepancy at \targetYear\ for \countryName.}
\end{figure}


% ============================================================================
\section{Values: IEA target vs SISEPUEDE before/after}
% ============================================================================

The following table reports the IEA target and the SISEPUEDE aggregate
(both in TJ) for each (balance, product) at the calibration target
year, before and after calibration. Rows include both sector totals
and their per-fuel components.

\input{\tableFile{values_table}}


% ============================================================================
\section{Calibration improvement}
% ============================================================================

The improvement table compares per-pair error before and after
calibration. The ratio column is colour-graded
(green: improvement, red: regression); the $\Delta$~(pp) column shows
the signed change in percentage points.

\input{\tableFile{improvement_table}}


% ============================================================================
\section{Before/after discrepancy plots}
% ============================================================================

\begin{figure}[H]
    \centering
    \begin{subfigure}{0.49\textwidth}
        \includegraphics[width=\linewidth]{\plotFile{before_after_discrepancy_ratio}}
        \caption{Ratio SSP/IEA.}
    \end{subfigure}\hfill
    \begin{subfigure}{0.49\textwidth}
        \includegraphics[width=\linewidth]{\plotFile{before_after_discrepancy_rel_error}}
        \caption{Relative error.}
    \end{subfigure}
    \caption{Discrepancy at \targetYear\ before vs.\ after calibration.
    Green bars indicate pairs where calibration moved closer to the
    IEA reference; red bars indicate regressions.}
\end{figure}


% ============================================================================
\section{Primary sector totals}
% ============================================================================

Phase 1 targets the row totals (INDUSTRY, TRANSPORT, RESIDENT,
COMMPUB, AGRICULT). The figures below show, for each primary sector,
the IEA series alongside the SISEPUEDE series before and after
calibration.

\begin{figure}[H]
    \centering
    \includegraphics[width=\linewidth]{\plotFile{primary_timeseries}}
    \caption{Primary sector totals --- IEA vs.\ SISEPUEDE before/after
    calibration, time series.}
\end{figure}

\begin{figure}[H]
    \centering
    \includegraphics[width=0.85\linewidth]{\plotFile{primary_bar}}
    \caption{Primary sector totals at \targetYear\ --- grouped bars.}
\end{figure}


% ============================================================================
\section{Fuel mix per sector}
% ============================================================================

Phase 2 targets the column shares within each sector. The blocks below
report the fuel composition (time series and target-year bars) for
each calibrated sector. Duplicate the block once per sector calibrated
for this country.

% ----------------------------------------------------------------------------
% TEMPLATE BLOCK — duplicate per sector, replacing SECTOR with the IEA
% balance code (e.g. INDUSTRY, TRANSPORT, RESIDENT, COMMPUB, AGRICULT).
% ----------------------------------------------------------------------------
\subsection{INDUSTRY}

\begin{figure}[H]
    \centering
    \includegraphics[width=\linewidth]{\plotFile{fuelmix_INDUSTRY_timeseries}}
    \caption{INDUSTRY fuel mix --- time series with ratio and relative-error
    diagnostic panels.}
\end{figure}

\begin{figure}[H]
    \centering
    \includegraphics[width=0.85\linewidth]{\plotFile{fuelmix_INDUSTRY_bar}}
    \caption{INDUSTRY fuel shares at \targetYear.}
\end{figure}

\newpage
\subsection{TRANSPORT}
\begin{figure}[H]
  \centering
  \includegraphics[width=\linewidth]{\plotFile{fuelmix_TRANSPORT_timeseries}}
  \caption{TRANSPORT fuel mix --- time series.}
\end{figure}
\begin{figure}[H]
  \centering
  \includegraphics[width=0.85\linewidth]{\plotFile{fuelmix_TRANSPORT_bar}}
  \caption{TRANSPORT fuel shares at \targetYear.}
\end{figure}


\newpage
\subsection{RESIDENT}
\begin{figure}[H]
  \centering
  \includegraphics[width=\linewidth]{\plotFile{fuelmix_RESIDENT_timeseries}}
  \caption{RESIDENT fuel mix --- time series.}
\end{figure}
\begin{figure}[H]
  \centering
  \includegraphics[width=0.85\linewidth]{\plotFile{fuelmix_RESIDENT_bar}}
  \caption{RESIDENT fuel shares at \targetYear.}
\end{figure}


\newpage
\subsection{COMMPUB}
\begin{figure}[H]
  \centering
  \includegraphics[width=\linewidth]{\plotFile{fuelmix_COMMPUB_timeseries}}
  \caption{COMMPUB fuel mix --- time series.}
\end{figure}
\begin{figure}[H]
  \centering
  \includegraphics[width=0.85\linewidth]{\plotFile{fuelmix_COMMPUB_bar}}
  \caption{COMMPUB fuel shares at \targetYear.}
\end{figure}


% ============================================================================
\section{Calibration knobs}
% ============================================================================

For each calibrated (balance, product) pair, the table below lists the
input variables that act as knobs for that pair, grouped by simplex
group ID. For fuel-mix targets (Phase 2), this means the
\texttt{frac\_*} fractions sitting in the same simplex constraint;
phase-1 scalar knobs appear under an ``Unconstrained'' sub-header. The
\textit{\% change} column is shaded yellow / orange / red when
$|\Delta| \geq 25 / 50 / 75\,\%$.

Each pair has its own file at
\texttt{\tablesDir/knobs\tagSuffix/knobs\_<balance>\_<product>.tex}; the
\textbackslash knobIfExists macro below pulls in only those that the
pipeline actually wrote, so unused pairs are silently skipped.

\newcommand{\knobIfExists}[2]{%
    \IfFileExists{\knobFile{#1}{#2}.tex}{%
        \subsection{\MakeUppercase{#1} $\times$ \MakeUppercase{#2}}%
        \input{\knobFile{#1}{#2}}%
    }{}%
}

% ── Sector totals (Phase 1 scalar groups) ───────────────────────────────────
\knobIfExists{industry}{industry}
\knobIfExists{transport}{transport}
\knobIfExists{resident}{resident}
\knobIfExists{commpub}{commpub}
\knobIfExists{agricult}{agricult}

% ── Fuel-mix fractions (Phase 2 simplex groups) ─────────────────────────────
% INDUSTRY × fuel
\knobIfExists{industry}{coal}
\knobIfExists{industry}{oil}
\knobIfExists{industry}{natgas}
\knobIfExists{industry}{electr}
\knobIfExists{industry}{biofuels}
\knobIfExists{industry}{heat}

% TRANSPORT × fuel
\knobIfExists{transport}{oil}
\knobIfExists{transport}{electr}
\knobIfExists{transport}{natgas}
\knobIfExists{transport}{biofuels}

% RESIDENT × fuel
\knobIfExists{resident}{coal}
\knobIfExists{resident}{oil}
\knobIfExists{resident}{natgas}
\knobIfExists{resident}{electr}
\knobIfExists{resident}{biofuels}
\knobIfExists{resident}{heat}

% COMMPUB × fuel
\knobIfExists{commpub}{coal}
\knobIfExists{commpub}{oil}
\knobIfExists{commpub}{natgas}
\knobIfExists{commpub}{electr}
\knobIfExists{commpub}{biofuels}
\knobIfExists{commpub}{heat}

% AGRICULT × fuel
\knobIfExists{agricult}{oil}
\knobIfExists{agricult}{electr}
\knobIfExists{agricult}{natgas}
\knobIfExists{agricult}{biofuels}

% Add additional \knobIfExists{<balance>}{<product>} lines for any other
% pairs the pipeline emits — the macro silently skips missing files.


\end{document}
